% SPDX-License-Identifier: Apache-2.0
\section{Output: FST through libfst}
\label{sec:output}

\subsection{Why not VCD}

The converter's output format has to hold what the database holds. VCD
does not: the measurement in the introduction, that eight of the
fifteen types in the corpus are absent from Vivado's own VCD, applies
to any VCD. FST~\cite{fst} has variable types for integers, reals,
times and strings, and it compresses far better.

\subsection{Why not a second FST writer}

FST has no standard. The nearest thing to a specification is the
source of \texttt{libfst}~\cite{libfst}, the library GTKWave, GHDL and
nvc all use, plus an unofficial write-up~\cite{hutt}. A second writer
would therefore be a second thing to maintain against a definition
that can move, with nothing to signal that it had drifted.

The converter calls \texttt{libfst} through cgo. The library is three
C source files and zlib, built by Bazel, so the bytes of the output
file come from the reference implementation. The same library reads
the output back in the test.

\subsection{The mapping}

FST has no record type and no array type, so the converter flattens:

\begin{itemize}
  \item a record becomes a scope with one variable per field,
        recursively;
  \item an array of anything but logic becomes a scope with one
        variable per element;
  \item a vector of logic stays one variable of that many bits;
  \item an integer is 32 bits of two's complement, a real is the
        double, and an enumeration, a character, a string or a
        physical value keeps its text.
\end{itemize}

Flattening was chosen so that a viewer can show, search and compare a
single field on its own. The aggregate can be read back from the
leaves; the reverse does not hold.

Two objects that share a handle are declared as one FST variable and
an alias of it, so a signal seen from several scopes is written once.

\subsection{Streaming}

The database stores each object's records where that object's storage
lies, and a waveform file wants every object's changes in ascending
time. The converter therefore replays the file once, keeping the
current value of every object and merging the arenas by record time.
The state is the handle space, 2.8~MB for the largest design here,
rather than the change history, which for that design is
18\,875\,466 entries.

\subsection{Measurements}

\begin{table}[t]
\caption{Conversion, measured on one machine}
\label{tab:cost}
\centering
\begin{tabular}{lrrrr}
\toprule
Design & \texttt{.wdb} & \texttt{.vcd} & \texttt{.fst} & Time \\
\midrule
\texttt{counter} & 14\,kB & 5.6\,kB & 1.6\,kB & 0.02\,s \\
\texttt{uart} & 45\,kB & 3.3\,kB & 2.3\,kB & 0.03\,s \\
PicoRV32 & 268\,kB & 287\,kB & 14\,kB & 0.25\,s \\
Potato & 197\,kB & 28\,kB & 38\,kB & 0.06\,s \\
SERV & 15.9\,MB & 29.6\,MB & 654\,kB & 8.8\,s \\
NEORV32 & 24.1\,MB & 17.2\,MB & 690\,kB & 71\,s \\
\bottomrule
\end{tabular}
\end{table}

Table~\ref{tab:cost} gives the cost. On the two long traces the FST is
45 and 25 times smaller than the VCD of the same run while holding
more. Potato is the instructive row: its FST is larger than its VCD,
because the VCD holds none of that design's integers, enumerations,
records or its two 32\,768 byte memories, and the FST holds all of
them.

Time is dominated by decoding a value and rendering its leaves for
every change, about 3.8~\textmu s per change. Memory is dominated by
holding the inflated pages of the whole database.
